\section{WorldFlow}
\label{sec:method}

\subsection{Problem Formulation}
\label{sec:problem}

Given a language instruction $c$, an observed RGB frame $I_0$,
and 3D query points
$\mathcal{Q}=\{\mathbf{p}_{0,i}\}_{i=1}^{N}$
with corresponding image coordinates
$\mathcal{S}_0=\{\mathbf{s}_{0,i}\}_{i=1}^{N}$,
\method{} jointly predicts future RGB frames $I_{1:T}$
and 3D flow $f_{1:T}$.
All 3D positions are expressed in the initial camera
coordinate system.
We define flow as the frame-to-frame displacement
of each query point:
\begin{equation}
f_{t,i}
=
\mathbf{p}_{t,i}-\mathbf{p}_{t-1,i},
\qquad t=1,\ldots,T.
\label{eq:flow-definition}
\end{equation}
The corresponding trajectories are recovered by accumulation:
\begin{equation}
\mathbf{p}_{t,i}
=
\mathbf{p}_{0,i}+\sum_{\tau=1}^{t}f_{\tau,i}.
\label{eq:trajectory-reconstruction}
\end{equation}
Our objective is to model the conditional joint distribution
\begin{equation}
p_\theta\left(
I_{1:T},f_{1:T}
\mid I_0,\mathcal{Q},\mathcal{S}_0,c
\right).
\label{eq:joint-distribution}
\end{equation}

\subsection{Architecture}
\label{sec:architecture}

Our central design introduces 3D flow tokens as an additional
generative modality for joint video-motion prediction.
Built on the Cosmos-Nano framework~\citep{nvidia2026cosmos3}, \method{} retains
the pretrained reasoner and extends the generator to
jointly synthesize future video and 3D flow
(Figure~\ref{fig:overview}).
The reasoner processes the language instruction, while
the video VAE encodes the observed frame into clean latents.
The flow tokenizer embeds noisy displacements together
with initial query-point geometry and aligns their
spatiotemporal positional coordinates with the video
token grid (Section~\ref{sec:tokenizer}).

Within the Mixture-of-Transformers (MoT) architecture,
video and flow tokens interact through full attention
and attend to the reasoner's causal context.
At inference, future video latents and flow are initialized
from noise and jointly denoised, conditioned on the clean
observed-frame latents, language context, and query geometry.
The VAE decodes the generated video latents into RGB frames,
while the predicted flow is rescaled and accumulated
to recover 3D trajectories.
This joint generation process produces explicit motion
alongside video without a separate tracking stage.

\begin{figure}[t]
    \centering
    \includegraphics[width=\linewidth]{figures/architecture.pdf}
    \caption{
        Overview of \method{}.
        Conditioned on text tokens $t_j$ through the pretrained
        reasoner, clean observed-frame latents $v_0$, and query-point geometry, the generator jointly denoises video and flow from noise $\epsilon_j$ and $\delta_k$ through shared multimodal attention.
        Generated video latents $v_j$ are decoded by the VAE,
        while generated 3D flow $f_t$ is rescaled and
        accumulated from the initial query positions
        to recover metric trajectories.
        The 2D flow overlay shows generated 3D trajectories projected onto generated video frames.
        The 3D flow overlay shows generated 3D trajectories overlaid on a scene point cloud reconstructed from
        the initial frame's ground-truth depth.
    }
    \label{fig:overview}
\end{figure}

\subsection{3D Flow Tokenization}
\label{sec:tokenizer}

We represent each query point's displacement at each future
frame as a continuous-valued token, augmented with initial
3D geometry and spatiotemporal positional encoding.

\paragraph{Normalization and sequence construction.}
To balance motion scales across domains, we normalize
displacements by a domain-specific scalar estimated
from training data and fixed throughout training:
\begin{equation}
\bar{f}_{t,i}=\frac{f_{t,i}}{a_d},
\qquad a_d>0.
\label{eq:motion-normalization}
\end{equation}
The normalized displacements are arranged in frame-major
order, yielding $TN$ three-dimensional vectors for
token construction.

\paragraph{Query point geometry encoder.}
We encode each query point's absolute initial position
and its position relative to the query-set centroid
using two parallel encoders.
Each applies coordinate normalization, fixed Fourier
features, and an MLP.
Their outputs are concatenated and fused:
\begin{equation}
\mathbf{e}^{\mathrm{3D}}_i
=
\mathrm{MLP}
\left(
\left[
E_{\mathrm{abs}}(\mathbf{p}_{0,i});
E_{\mathrm{rel}}(\mathbf{p}_{0,i}-\boldsymbol{\mu})
\right]
\right),
\label{eq:point-pe}
\end{equation}
where $E_{\mathrm{abs}}$ and $E_{\mathrm{rel}}$ denote
the two encoders, and $\boldsymbol{\mu}$ is the query-set
centroid.
The geometry embedding is shared across all future
tokens of the same query point.

\paragraph{Flow token construction.}
A linear projection maps noisy normalized displacements
into the generator's hidden space, where they are combined
with geometry and conditioning embeddings:
\begin{equation}
\mathbf{h}_{t,i}^{(\sigma)}
=
\mathrm{Linear}(\bar{f}_{t,i}^{(\sigma)})
+
\mathbf{e}^{\mathrm{3D}}_i
+
\mathbf{e}_{\mathrm{cond}}(\sigma).
\label{eq:flow-embedding}
\end{equation}
Here, $\sigma$ denotes the denoising noise level,
distinct from the physical frame index $t$.
The conditioning embedding combines a learned
flow-modality embedding with the noise-level embedding.

\paragraph{Spatiotemporal alignment.}
We assign flow tokens rotary positional coordinates
using the same temporal and spatial conventions as
video tokens.
Each flow token is anchored to its future frame
and the query point's initial image patch:
\begin{equation}
\mathbf{r}_{t,i}
=
\left(
\frac{t}{q_t},
\left\lfloor\frac{v_{0,i}}{q_s p_s}\right\rfloor,
\left\lfloor\frac{u_{0,i}}{q_s p_s}\right\rfloor
\right),
\label{eq:flow-rope-coordinates}
\end{equation}
where $(u_{0,i},v_{0,i})=\mathbf{s}_{0,i}$ denotes
the initial image coordinates, $q_t$ and $q_s$ are
the temporal and spatial VAE compression factors,
and $p_s$ is the latent patch size.
These coordinates determine the rotary transformations
of attention queries and keys:
\begin{equation}
\widetilde{\mathbf{q}}_{t,i}
=
R(\mathbf{r}_{t,i})\mathbf{q}_{t,i},
\qquad
\widetilde{\mathbf{k}}_{t,i}
=
R(\mathbf{r}_{t,i})\mathbf{k}_{t,i}.
\label{eq:flow-rope}
\end{equation}
The shared positional convention encodes relative
spatiotemporal relationships between flow and video
tokens within multimodal attention.

\subsection{Training Objectives}
\label{sec:training}

\subsubsection{Joint flow matching}
We jointly train video and 3D flow generation using
a weighted rectified-flow objective:
\begin{equation}
\mathcal{L}_{\mathrm{FM}}
=
\lambda_{\mathrm{video}}\mathcal{L}_{\mathrm{video}}
+
\lambda_{\mathrm{flow}}\mathcal{L}_{\mathrm{flow}},
\label{eq:joint-objective}
\end{equation}
where $\mathcal{L}_{\mathrm{video}}$ and
$\mathcal{L}_{\mathrm{flow}}$ supervise vector-field
prediction for video latents and normalized 3D flow,
respectively.
Both modalities share a per-sample noise level and
are jointly denoised.
Losses are computed separately for each modality,
excluding conditioning positions from supervision.

\subsubsection{Auxiliary motion objectives}
To complement vector-field supervision with explicit
constraints on predicted motion, we apply auxiliary
objectives to clean normalized flow reconstructed
from the predicted vector field.
These objectives penalize frame-to-frame displacement
error, cumulative displacement direction error,
and relative endpoint error, encouraging both local
accuracy and trajectory-level consistency.

Direction errors are weighted by ground-truth displacement
magnitude, while points with small net displacement
are excluded from endpoint supervision.
All auxiliary terms operate in normalized coordinates
and use noise-dependent weighting to emphasize
lower-noise samples.
The complete objective is
\begin{equation}
\mathcal{L}
=
\mathcal{L}_{\mathrm{FM}}
+
\lambda_{\mathrm{dir}}\mathcal{L}_{\mathrm{dir}}
+
\lambda_{\mathrm{vmag}}\mathcal{L}_{\mathrm{vmag}}
+
\lambda_{\mathrm{end}}\mathcal{L}_{\mathrm{end}},
\label{eq:complete-objective}
\end{equation}
where $\mathcal{L}_{\mathrm{dir}}$,
$\mathcal{L}_{\mathrm{vmag}}$, and
$\mathcal{L}_{\mathrm{end}}$ denote the direction,
displacement, and endpoint losses, respectively.
Detailed formulations and derivations are provided
in the supplementary material, and their contribution
is evaluated through ablation experiments.